\documentclass[../../main.tex]{subfiles}
\begin{document}
\section{超定偏微分方程组}\label{section:ch19-overdetermined}

我们在第 9 章中考虑的偏微分方程都是关于一个未知函数的单个方程。在某些应用中，需要考虑\emph{超定的}偏微分方程组，即方程个数多于未知函数个数的方程组。一般而言，超定方程组只有在满足某些相容条件时才有解。对某些一阶方程组，相容条件可解释为关于某分布对合性的命题，而 Frobenius 定理可用来证明解的局部存在唯一性。

首先考虑某些线性方程组。设 $W$ 是 $\mathbb{R}^n$ 的一个开子集，$m$ 是一个不超过 $n$ 的正整数。考虑关于单个未知函数 $u\in C^\infty(W)$ 的、由 $m$ 个线性偏微分方程组成的下列方程组：
\begin{align}
a^1_1(x)\frac{\partial u}{\partial x^1}(x)+\cdots+a^n_1(x)\frac{\partial u}{\partial x^n}(x)&=f_1(x),\nonumber\\
\vdots\qquad\qquad&\\
a^1_m(x)\frac{\partial u}{\partial x^1}(x)+\cdots+a^n_m(x)\frac{\partial u}{\partial x^n}(x)&=f_m(x),\nonumber
\end{align}
其中 $\bigl(a^j_i\bigr)$ 是一个 $n\times m$ 的光滑实值函数矩阵，$f_1,\dots,f_m$ 是 $W$ 上的光滑实值函数。$m=1$ 的情形已由定理\cref{theorem:ch09-51} 涵盖，故这里的讨论主要在 $m>1$ 时有用。

若令 $A_i$ 表示向量场 $a^j_i\,\partial/\partial x^j$，则方程组可更简洁地写成 $A_1u=f_1,\dots,A_mu=f_m$。为避免冗余或退化的方程组，我们假设矩阵 $\bigl(a^j_i\bigr)$ 在 $W$ 的每一点秩为 $m$，等价地，向量场 $A_1,\dots,A_m$ 线性无关。

下列定理是定理\cref{theorem:ch09-51} 在超定情形的类比。

\begin{theorem}\label{theorem:ch19-超定线性}
设 $W\subseteq\mathbb{R}^n$ 是开子集，$m$ 是满足 $1\leqslant m\leqslant n$ 的整数。给定一个嵌入的余维-$m$ 子流形 $S\subseteq W$、$W$ 上一组线性无关的 $m$ 个光滑向量场 $(A_1,\dots,A_m)$，其张成空间在每个 $p\in S$ 处与 $T_pS$ 互补，以及函数 $f_1,\dots,f_m\in C^\infty(W)$。又设存在光滑函数 $c^k_{ij}\in C^\infty(W)$（$i,j,k=1,\dots,m$），满足下列相容条件：
\begin{align}
[A_i,A_j]&=c^k_{ij}A_k,\\
A_i f_j-A_j f_i&=c^k_{ij}f_k.
\end{align}
（在这些式中，$k$ 从 $1$ 到 $m$ 隐式求和。）则对每个 $p\in S$，存在 $p$ 的一个邻域 $U$，使得对每个 $\varphi\in C^\infty(S\cap U)$，都存在 $u\in C^\infty(U)$ 满足 $A_i u=f_i$（$i=1,\dots,m$）以及初始条件 $u|_{S\cap U}=\varphi$。
\end{theorem}

\begin{proof}
向量场 $(A_1,\dots,A_m)$ 张成 $W$ 上一个光滑的秩-$m$ 分布。由相容条件与引理\cref{lemma:ch19-局部标架对合判别}，该分布对合。于是由 Frobenius 定理，在 $p$ 的一个邻域内可取平坦坐标。在这样的坐标下，$A_i$ 张成的分布的叶由坐标 $(v,w)=(v^1,\dots,v^m,w^1,\dots,w^{n-m})$ 中 $w$ 为常数给出。

令 $(\alpha^1,\dots,\alpha^m,\beta^1,\dots,\beta^{n-m})$ 表示对偶余标架，并定义光滑 $1$-形式 $\omega\in\Omega^1(U)$ 为 $\omega=f^k\alpha^k$（从 $1$ 到 $m$ 隐式求和）。方程组 $A_i u=f_i$ 等价于 $\mathrm{d}u(A_i)=\omega(A_i)$（$i=1,\dots,m$），这等价于说 $\mathrm{d}u-\omega$ 拉回到每个叶 $H_w$ 上等于零。

利用外导数公式\eqref{align:ch14-d1form-invariant} 以及相容条件，我们得到对每个 $i,j,k=1,\dots,m$，
\begin{align*}
\mathrm{d}\alpha^k(A_i,A_j)=A_i\bigl(\alpha^k(A_j)\bigr)-A_j\bigl(\alpha^k(A_i)\bigr)-\alpha^k\bigl([A_i,A_j]\bigr)=c^k_{ij}.
\end{align*}
于是由相容条件的第二式，
\begin{align*}
\mathrm{d}\omega(A_i,A_j)=(A_i f_k)\delta^k_j-(A_j f_k)\delta^k_i-f_k c^k_{ij}=A_i f_j-A_j f_i-f_k c^k_{ij}=0.
\end{align*}
由于 $(A_1,\dots,A_m)$ 在每个积分流形 $H_w$ 上限制为一组标架，这表明 $\omega$ 拉回到每个 $H_w$ 上是闭的。

给定 $\varphi\in C^\infty(U\cap S)$，令 $u=u_0+u_1$，其中 $u_0,u_1\in C^\infty(U)$ 定义为
\begin{align*}
u_0(v,w)=\varphi(0,w),\qquad
u_1(v,w)=\int_0^1 \omega^k(tv,w)\,v^k\,\mathrm{d}t,
\end{align*}
其中 $\omega^k\,\mathrm{d}v^k$ 是 $\omega$ 的坐标表示。回忆平坦坐标图按定义是立方体，因而星形，故积分对所有 $(v,w)\in U$ 都良定义；在积分号下求导表明 $u_1$ 是 $(v,w)$ 的光滑函数。由于 $u_0|_{S\cap U}=\varphi$、$u_1|_{S\cap U}=0$，可知 $u$ 满足初始条件。

函数 $u_0$ 满足 $A_1u_0=\cdots=A_mu_0=0$，因为它与 $v$ 坐标无关。另一方面，对每个固定的 $w$，$u_1$ 是 $H_w$ 上由公式\eqref{align:ch11-势函数} 给出的、对应于 $\omega|_{H_w}$ 的势函数。定理\cref{theorem:ch11-49} 的证明表明该势函数满足 $\mathrm{d}u_1|_{H_w}=\omega|_{H_w}$，因而 $u=u_0+u_1$ 满足方程组。
\end{proof}

接下来考虑一个（可能非线性的）一阶方程组。设我们要求解
\begin{align}
\frac{\partial u}{\partial x}=\alpha(x,y,u),\qquad \frac{\partial u}{\partial y}=\beta(x,y,u),
\end{align}
其中 $\alpha$ 与 $\beta$ 是定义在某个开子集 $W\subseteq\mathbb{R}^3$ 上的光滑实值函数。这是一个关于（可能非线性的）一阶 PDE 的超定方程组。（事实上，对两个自变量的未知函数，几乎任意一对光滑一阶偏微分方程组至少在局部都可写成这种形式，只需解出 $\partial u/\partial x$ 与 $\partial u/\partial y$ 即可。能否在原则上这样做是一个由隐函数定理完全回答的问题；而能否在实际上这样做则取决于具体方程以及你的技巧。）

为确定 $\alpha,\beta$ 使方程组有解所必须满足的相容条件，设 $u$ 是 $\mathbb{R}^2$ 的某个开子集上的光滑解。由于 $\partial^2u/\partial x\partial y=\partial^2u/\partial y\partial x$，方程组蕴含
\begin{align*}
\frac{\partial}{\partial y}\alpha\bigl(x,y,u(x,y)\bigr)=\frac{\partial}{\partial x}\beta\bigl(x,y,u(x,y)\bigr),
\end{align*}
从而由链式法则
\begin{align}
\frac{\partial\alpha}{\partial y}+\beta\frac{\partial\alpha}{\partial z}=\frac{\partial\beta}{\partial x}+\alpha\frac{\partial\beta}{\partial z}.
\end{align}
只要存在光滑解 $u$ 满足 $u(x,y)=z$，这就在点 $(x,y,z)\in W$ 处成立。特别地，上式是方程组在每个点 $(x_0,y_0)$ 的一个邻域内、且初值 $u(x_0,y_0)=z_0$ 可任意指定时有解的必要条件。利用 Frobenius 定理，可证明该条件也是充分的。

\begin{proposition}\label{proposition:ch19-超定非线性}
设 $\alpha,\beta$ 是定义在某个开子集 $W\subseteq\mathbb{R}^3$ 上的光滑实值函数，并在该处满足上式。对每个 $(x_0,y_0,z_0)\in W$，存在 $(x_0,y_0)$ 在 $\mathbb{R}^2$ 中的一个邻域 $U$ 以及唯一的光滑函数 $u\colon U\to\mathbb{R}$，满足上述方程组以及 $u(x_0,y_0)=z_0$。
\end{proposition}

\begin{proof}
证明的想法是：该方程组用 $u$ 的值确定了 $u$ 的偏导数，因而用图像上点的坐标确定了图像在每一点处的切平面。这一族切平面在 $W$ 上定义了一个光滑的秩-$2$ 分布（见图\cref{figure:ch19-graph}），而上述相容条件等价于该分布的对合性条件。

若在开子集 $U\subseteq\mathbb{R}^2$ 上有解 $u$，则由 $F\colon U\to W$，$F(x,y)=(x,y,u(x,y))$ 给出的映射满足方程组当且仅当
\begin{align*}
\mathrm{d}F\left(\frac{\partial}{\partial x}\bigg|_{(x,y)}\right)&=\frac{\partial}{\partial x}\bigg|_p+\alpha\bigl(x,y,u(x,y)\bigr)\frac{\partial}{\partial z}\bigg|_p,\\
\mathrm{d}F\left(\frac{\partial}{\partial y}\bigg|_{(x,y)}\right)&=\frac{\partial}{\partial y}\bigg|_p+\beta\bigl(x,y,u(x,y)\bigr)\frac{\partial}{\partial z}\bigg|_p.
\end{align*}
令 $X,Y$ 为 $W$ 上的向量场
\begin{align}
X=\frac{\partial}{\partial x}+\alpha(x,y,z)\frac{\partial}{\partial z},\qquad
Y=\frac{\partial}{\partial y}+\beta(x,y,z)\frac{\partial}{\partial z},
\end{align}
并令 $D$ 为 $W$ 上由 $X,Y$ 张成的分布。由于上式恰说明 $F$ 的切空间等于 $D$ 的叶的切空间，$F$ 是积分流形之间的映射。由计算 $[X,Y]$ 可知相容条件等价于 $D$ 对合；于是由 Frobenius 定理，过 $(x_0,y_0,z_0)$ 有 $D$ 的一个积分流形。在适当的坐标下该积分流形是一个图像 $z=u(x,y)$，它给出所求的解；唯一性由积分流形的局部结构得到。具体解出 $w=\text{常数}$ 中的 $z=u(x,y)$ 即可。一些例子见习题 19-13。
\end{proof}

\begin{figure}[ht]
\centering
\includegraphics[width=0.8\textwidth]{fig-19-10.png}
\caption{通过寻找图像求解 $u$}\label{figure:ch19-graph}
\end{figure}

这个结果可直接推广到更高维。该定理的一般叙述稍为复杂，但在具体例子中验证必要条件通常只涉及计算混合偏导数并应用链式法则。

\begin{proposition}\label{proposition:ch19-高维超定}
设 $W$ 是 $\mathbb{R}^n\times\mathbb{R}^m$ 的一个开子集，$\alpha=\bigl(\alpha^i_j\bigr)\colon W\to M(m\times n,\mathbb{R})$ 是一个光滑矩阵值函数，满足
\begin{align*}
\frac{\partial\alpha^i_j}{\partial x^k}+\alpha^l_k\frac{\partial\alpha^i_j}{\partial z^l}
=\frac{\partial\alpha^i_k}{\partial x^j}+\alpha^l_j\frac{\partial\alpha^i_k}{\partial z^l}
\end{align*}
对所有 $i,j,k$ 成立，其中把 $\mathbb{R}^n\times\mathbb{R}^m$ 中的点记作 $(x,z)=(x^1,\dots,x^n,z^1,\dots,z^m)$。对任意 $(x_0,z_0)\in W$，存在 $x_0$ 在 $\mathbb{R}^n$ 中的一个邻域 $U$ 以及唯一的光滑函数 $u\colon U\to\mathbb{R}^m$，使得 $u(x_0)=z_0$，且 $u$ 的 Jacobi 矩阵满足
\begin{align*}
\frac{\partial u^i}{\partial x^j}(x^1,\dots,x^n)=\alpha^i_j\bigl(x^1,\dots,x^n,u^1(x),\dots,u^m(x)\bigr).
\end{align*}
\end{proposition}

\begin{exercise}\label{exercise:ch19-30}
证明命题\cref{proposition:ch19-高维超定}。
\end{exercise}
\end{document}
